\indent ಉದಾಹರಣೆಗೆ,\par

\begin{kbvderiv}
\vspace{4pt plus 3pt minus 1pt}
\begin{center}\begin{adjustbox}{max width=\textwidth}
{\renewcommand{\arraystretch}{1.22}\setlength{\tabcolsep}{4pt}
\begin{tabular}{l c l}
(ನಾಮ ಪ್ರಕೃತಿ) + (ಪ್ರತ್ಯಯ) & {\color{kbvgrey}$=$} & \kbvresult{(ನಾಮಪದ)} \\
ಲಕ್ಷ್ಮಣ + ಉ & {\color{kbvgrey}$=$} & \kbvresult{ಲಕ್ಷ್ಮಣನು} \\
ರಾಮ + ಅ & {\color{kbvgrey}$=$} & \kbvresult{ರಾಮನ} \\
ಸೀತೆ + ಅನ್ನು & {\color{kbvgrey}$=$} & \kbvresult{ಸೀತೆಯನ್ನು} \\
ಅರಣ್ಯ + ಅಲ್ಲಿ & {\color{kbvgrey}$=$} & \kbvresult{ಅರಣ್ಯದಲ್ಲಿ} \\
\end{tabular}}
\end{adjustbox}\end{center}
\vspace{4pt plus 3pt minus 1pt}
\end{kbvderiv}

\kbvexercise{ಅಭ್ಯಾಸ ಪತ್ರಿಕೆ ೩}

\vspace{2.5pt plus 3pt minus 0.5pt}\noindent\hspace*{1.2em}\kbvnum{(1)}\hspace{0.55em}ನಾಮಪದ ಎಂದರೇನು ? ಉದಾಹರಣೆಗಳ ಸಹಿತ ತಿಳಿಸಿ.\par

\vspace{2.5pt plus 3pt minus 0.5pt}\noindent\hspace*{1.2em}\kbvnum{(2)}\hspace{0.55em}ನಾಮ ಪ್ರಕೃತಿಗೂ ನಾಮಪದಕ್ಕೂ ಇರುವ ವ್ಯತ್ಯಾಸವೇನು?\par

\vspace{2.5pt plus 3pt minus 0.5pt}\noindent\hspace*{1.2em}\kbvnum{(3)}\hspace{0.55em}ಕೆಳಗೆ ಕೊಟ್ಟಿರುವ ವಾಕ್ಯಗಳಲ್ಲಿ ನಾಮಪದಗಳನ್ನು ಗುರ್ತಿಸಿ ಮತ್ತು ಪ್ರತಿಯೊಂದು ನಾಮಪದವನ್ನೂ ಬಿಡಿಸಿ ಪ್ರಕೃತಿ, ಪ್ರತ್ಯಯ ತಿಳಿಸಿ\par

\vspace{2.5pt plus 3pt minus 0.5pt}\noindent\hspace*{1.2em}\kbvnum{(೧)}\hspace{0.55em}ಭೀಮನು ಕೀಚಕನನ್ನು ಕೊಂದನು.\par

\vspace{2.5pt plus 3pt minus 0.5pt}\noindent\hspace*{1.2em}\kbvnum{(೨)}\hspace{0.55em}ಕೋತಿಯು ಶಾಲಿನಿಯಿಂದ ಹಣ್ಣನ್ನು ಕಸಿಯಿತು.\par

\vspace{2.5pt plus 3pt minus 0.5pt}\noindent\hspace*{1.2em}\kbvnum{(೩)}\hspace{0.55em}ಕರ್ಣನು ಸರ್ಪಾಸ್ತ್ರವನ್ನು ಪ್ರಯೋಗಿಸಿದನು.\par

\vspace{2.5pt plus 3pt minus 0.5pt}\noindent\hspace*{1.2em}\kbvnum{(೪)}\hspace{0.55em}ಆಮೆಯು ನೀರಿನಲ್ಲಿ ವಾಸಿಸುತ್ತದೆ.\par

\vspace{2.5pt plus 3pt minus 0.5pt}\noindent\hspace*{1.2em}\kbvnum{(೫)}\hspace{0.55em}ಬೇಟೆಗಾರನು ಕಾಡಿಗೆ ಹೋದನು.\par

\kbvchapter{ಪಾಠ ೪}{ವಾಕ್ಯರಚನೆ}

\indent ವಸ್ತು ಸ್ಥಿತಿಯ ಅರಿವುಂಟುಮಾಡುವ ಶಬ್ದಗಳೇ ಪದಗಳು. ಪದಗಳ ರಚನೆಗೆ ಅಕ್ಷರಗಳು ಅತ್ಯಗತ್ಯ. ರೂಢಿಯಲ್ಲಿರುವ ಅರ್ಥಬರುವಂತೆ ಅಕ್ಷರಗಳನ್ನು ಕೂಡಿಸಿದಾಗಲೇ ಪದ ಉತ್ಪತ್ತಿಯಾಗುವುದು.\par

\indent ಹಲವು ಪದಗಳು ಸೇರಿ ವಾಕ್ಯವಾಗುವುದು. ವಾಕ್ಯವು ಅರ್ಥವತ್ತಾಗಿರ ಬೇಕಾದರೆ, ಆ ವಾಕ್ಯದಲ್ಲಿನ ಪದಗಳಲ್ಲಿ ಸೂಕ್ತವಾದ ಹೊಂದಾಣಿಕೆ ಇರಬೇಕು.\par

\indent ಉದಾಹರಣೆಗೆ ಕತ್ತೆಯು, ಭಾರವನ್ನು, ಹೊರುತ್ತದೆ ಎಂಬ ಪದ ಗಳನ್ನು ಗಮನಿಸೋಣ. ಇಲ್ಲಿ “ಕತ್ತೆಯು” ಎಂದಷ್ಟೇ ಹೇಳಿದರೆ ಮಾತಿನ ಅರ್ಥ ತಿಳಿದುಬರುವುದಿಲ್ಲ. ಹೀಗೆ, ಮೇಲೆ ತಿಳಿಸಿದ ಯಾವುದೊಂದು ಪದವನ್ನು\kbvcontinues
